\chapter[Referencial Teórico]{Referencial Teórico}
Este capítulo apresenta a revisão de literatura que fundamenta o trabalho, reunindo conceitos, definições e estudos relevantes sobre quatro eixos: (i) esporte comunitário e espaço público; (ii) \textit{Computer-Supported Cooperative Work} (CSCW); (iii) experiências ``\textit{phygital}'', mapas e descoberta local; e (iv) cidades inteligentes. O objetivo é estabelecer a base conceitual necessária à compreensão do tema e dos desdobramentos discutidos nos capítulos seguintes, conforme a orientação de que o referencial teórico corresponda a uma análise dos trabalhos relevantes encontrados na pesquisa bibliográfica.

\section{Esporte comunitário e espaço público}
Neste trabalho, entende-se \textit{esporte comunitário} como a prática física realizada no cotidiano da cidade, de forma espontânea ou organizada, em espaços públicos como praças, parques e quadras. A literatura aponta que essas práticas vão além do lazer: contribuem para a saúde, fortalecem vínculos sociais, produzem capital social e qualificam o uso do espaço comum \cite{coalter2007}. Diretrizes de saúde urbana indicam que ambientes acessíveis, seguros e próximos, como calçadas adequadas, áreas verdes e equipamentos esportivos, favorecem níveis mais altos de atividade física quando combinados a oportunidades locais \cite{who2018gappa}. No campo do urbanismo, documentos recentes defendem que espaços públicos de qualidade promovem integração social, convivência e participação cívica, e por isso devem ser planejados, distribuídos e mantidos de modo inclusivo \cite{unhabitat2015publicspace}. Em síntese, o esporte comunitário pode ser lido como um arranjo de bem estar e coesão que acontece na e pela cidade, dependendo de equipamentos esportivos legíveis, cuidados e efetivamente disponíveis à população.

\section{CSCW (Computer Supported Cooperative Work)}
Neste trabalho, CSCW é entendido como o campo interdisciplinar que estuda e projeta tecnologias para apoiar a cooperação entre pessoas. A ênfase recai em tornar o trabalho em grupo mais coordenado e previsível, com menos desencontros e incertezas. Para isso, a literatura descreve conceitos e mecanismos que estruturam a colaboração, como o trabalho de articulação e a consciência situacional, além de apontar desafios recorrentes no desenvolvimento de sistemas colaborativos \cite{schmidt1992taking,dourish1992awareness,olson2000distance,grudin1994eight}.
    \subsection{Trabalho de articulação (\textit{articulation work})}
    O trabalho de articulação é o esforço necessário para combinar tarefas, recursos e papéis para que uma atividade coletiva aconteça do início ao fim. Ele envolve definir quem faz o que, quando e com quais dependências, além de lidar com imprevistos e ajustes de última hora \cite{schmidt1992taking}. Em contextos colaborativos, esse esforço se torna explícito por meio de regras simples e artefatos que ajudam a organizar o fluxo da atividade.
    \subsection{Consciência situacional (\textit{awareness})}
    Consciência situacional é a visibilidade contínua sobre quem participa, o que já foi feito, o que falta fazer e quando algo muda. Essa noção permite que as pessoas sincronizem ações, antecipem conflitos de agenda e respondam a alterações de contexto sem depender de trocas individuais constantes \cite{dourish1992awareness}. Em ambientes digitais, a consciência situacional costuma aparecer em sinais claros de estado, confirmações, avisos e indicadores de presença.
    \subsection{Artefatos e mecanismos de coordenação}
    Sistemas colaborativos eficazes oferecem artefatos que tornam a coordenação tangível e menos custosa para o grupo. Entre os mecanismos frequentes estão agendas compartilhadas, listas de presença, distribuição de papéis, regras de entrada e saída, estados compartilhados e lembretes que chegam no momento oportuno. Esses elementos ajudam a reduzir a assimetria de informação e a necessidade de renegociação constante \cite{olson2000distance}.
    \subsection{Desafios clássicos do \textit{groupware}}
    A literatura identifica obstáculos recorrentes na adoção e no uso de tecnologias colaborativas. Entre eles estão a assimetria de benefícios entre quem cadastra informações e quem as consome, o encaixe precário com rotinas reais, o custo adicional de registro e manutenção e a fragmentação de informações em múltiplos canais \cite{grudin1994eight}. Estudos também mostram que a coordenação se torna mais difícil com distância, fuso horário e atrasos de comunicação, o que exige feedbacks oportunos e estados compartilhados bem definidos \cite{olson2000distance}.

\section{Phygital}
Neste trabalho, usa se o termo phygital para descrever experiências em que camadas digitais e físicas se entrelaçam no cotidiano da cidade. Com a difusão dos smartphones, tornou se comum decidir no ambiente digital o que será feito no ambiente físico, com base em informações de contexto como localização, horário, disponibilidade e avaliações. Essa perspectiva dialoga com estudos sobre espaços híbridos, nos quais mobilidade e conectividade alteram a forma de sociabilidade urbana \cite{desouzaesilva2006hybrid}, e com a literatura de serviços baseados em localização, que sustenta funcionalidades como busca por pontos próximos, rotas e filtros de distância e tempo \cite{raper2007lbs}. Também se relaciona à noção de geografia voluntária, em que cidadãos registram lugares e eventos, alimentando mapas vivos de oportunidades locais \cite{goodchild2007citizens}.
    \subsection{Espaços híbridos (\textit{hybrid spaces})}
    A ideia de espaços híbridos descreve contextos em que redes digitais acompanham as pessoas em movimento e se combinam às situações presenciais. Nessas condições, o acesso a informação contextual deixa de ser separado do lugar e do momento e passa a compor a própria experiência urbana, influenciando encontros, decisões e rotas \cite{desouzaesilva2006hybrid}. O resultado é uma cidade na qual ver, decidir e agir caminham juntos em um mesmo fluxo.
    \subsection{Serviços baseados em localização (\textit{location based services})}
    Serviços baseados em localização são sistemas que usam posição geográfica para entregar informação relevante ao usuário. Entre seus componentes típicos estão mecanismos de posicionamento e geocodificação, busca por proximidade, traçado de rotas e filtros temporais para selecionar oportunidades por dia e hora \cite{raper2007lbs}. Em aplicações do cotidiano, esses serviços permitem responder a perguntas práticas como onde ir agora, como chegar e o que está acontecendo por perto.
    \subsection{Dados voluntários e mapas vivos (\textit{volunteered geographic information})}
    A noção de geografia voluntária descreve a produção de dados espaciais por cidadãos, que atuam como sensores ao registrar lugares, eventos, horários e condições do ambiente. Esses registros alimentam mapas vivos que refletem o que acontece no território em tempo quase real, ampliando a descoberta hiperlocal e a previsibilidade de encontros \cite{goodchild2007citizens}. Ao mesmo tempo, a literatura aponta cuidados com qualidade, validação e privacidade desses dados.
    \subsection{Descoberta hiperlocal e online to offline}
    Com a combinação de espaços híbridos, serviços baseados em localização e dados voluntários, consolidou se um padrão de uso em que as pessoas escolhem no digital o que farão no físico. A descoberta hiperlocal depende de sinais simples e atualizados, como localização, horário e disponibilidade, que reduzem a incerteza antes do deslocamento. Esse movimento de online para offline aproxima a decisão do momento da ação e torna mais eficiente a ocupação de lugares e a participação em atividades presenciais.

\section{Cidades inteligentes}

Neste trabalho, entende se cidade inteligente como uma abordagem que combina infraestrutura urbana, capital humano e social e tecnologias digitais para melhorar serviços públicos, sustentabilidade e qualidade de vida. Uma definição operacional amplamente utilizada indica que a inteligência urbana emerge do investimento coordenado em pessoas e em tecnologias de informação para promover desenvolvimento sustentável e bem estar \cite{caragliu2011smart}.

No eixo de mensuração, normas e recomendações internacionais propõem conjuntos de indicadores para acompanhar o desempenho urbano em domínios como mobilidade, saúde, segurança e lazer. Entre essas referências destacam se a série ISO 37122, dedicada a indicadores de cidades inteligentes em complemento à ISO 37120, e as recomendações do ITU T para indicadores de cidades inteligentes e sustentáveis \cite{iso37122,itutY4903}. Leituras críticas observam que a instrumentação e a análise de dados em tempo quase real criam novas capacidades de gestão, mas exigem governança clara de dados, transparência e participação para evitar dependências tecnológicas e vieses tecnocráticos \cite{kitchin2014realtime}. Assim, o debate contemporâneo privilegia uma visão sociotécnica em que dados e plataformas são meios para ampliar bem estar, inclusão e uso eficiente dos ativos públicos, e não fins em si.

\bigskip
\noindent\textbf{Síntese.} Os quatro eixos aqui revisados fornecem a base conceitual necessária para compreender o fenômeno do esporte comunitário na cidade, os princípios de coordenação mediados por tecnologia, as dinâmicas híbridas entre camadas digitais e físicas no espaço urbano e os marcos de mensuração e governança associados a cidades inteligentes. Nos capítulos seguintes, tais fundamentos amparam as escolhas metodológicas e as análises conduzidas neste trabalho.